\chapter{Software Engineer}\label{ch:in:software-engineer}

In February 2026 Virtu Financial, one of the largest listed electronic market makers, had about 1\,027 employees in
thirteen countries. Five months later Microsoft reported 77\,000 people in product research and development alone,
seventy-five times Virtu's whole staff. A software engineer at a trading firm is one of a few hundred engineers at most,
a few steps from the profit and loss, and on call while markets are open; at a technology company the same engineer is
one of tens of thousands, many steps from revenue, with a product that runs around the clock. This chapter describes the
engineering jobs at trading firms, shows how they differ from the same jobs in technology, and compares the pay the filings
and the survey record for them.

\begin{center}\footnotesize
\begin{tabular}{@{}p{2.4cm}p{3.3cm}p{3.3cm}p{3.3cm}@{}}
\toprule
\multicolumn{4}{@{}l}{\textbf{Role cards: three engineering specialisms}}\\
\midrule
 & \omterm{def:in:software-engineer:lowlat}{low-latency engineer} & \omterm{def:in:software-engineer:fpga}{FPGA engineer} & \omterm{def:in:software-engineer:sre}{site reliability engineer}\\
\midrule
works on & the software path from market data to order & logic in programmable hardware on that path & the production systems and their operation\\
units & microseconds and their tails & nanoseconds and clock cycles & availability, incidents, time to recover\\
reports to & head of trading technology & head of hardware & head of infrastructure\\
codes in filings & 15-1252 & 15-1252, 17-2061 & 15-1252, 15-1299.08\\
taught in & Book~13, ch.~1; Book~14, ch.~1 & Book~14, ch.~6--7 & Book~15, ch.~28; Book~16, ch.~21\\
\bottomrule
\end{tabular}
\end{center}

\section{Core, low-latency, infrastructure, hardware and network engineering}

A trading firm's engineers divide by what they build. \emph{Core} engineers build the trading system's logic and the
tools around it: the strategy engine, order management, risk checks, the configuration that traders change, and the
research engineering of chapter~19. The rest specialise by the layer they own.

\begin{definition}[Low-latency engineer]\label{def:in:software-engineer:lowlat}
A \emph{low-latency engineer}\index{low-latency engineer} is a software engineer whose work is to shorten and stabilise the
time between an event in the market and the firm's response to it: the code, memory layout, operating-system settings
and network stack on that path, measured in microseconds and in the tail of their distribution.
\end{definition}

The \omterm{def:in:software-engineer:lowlat}{low-latency engineer}'s craft is the subject of Book~13: where latency comes from (chapter~1), how processors, caches
and the operating system add to it, and the kernel bypass that removes the kernel from the network path (chapter~13).
The work is measured, not argued: a change ships when its latency histogram, replayed on the same data, is better where
it matters, usually in the tail.

\begin{definition}[FPGA engineer]\label{def:in:software-engineer:fpga}
An \emph{FPGA engineer}\index{FPGA engineer} is an engineer who designs, verifies and maintains logic for
field-programmable gate arrays: feed handlers, order builders and risk checks that run in hardware, where the time from a
packet's arrival to a response is counted in clock cycles.
\end{definition}

Book~14 covers the hardware (chapters~6 and~7) and the network around it: switches, network cards, time
synchronisation, colocation. The \omterm{def:in:software-engineer:fpga}{FPGA engineer} writes a hardware description language, verifies it in simulation far
more thoroughly than software is tested (a hardware bug cannot be patched in the middle of a trading day), and works with
the network engineers who own the cables, switches and the colocated racks. It is the smallest of the engineering
specialisms and the one most specific to trading.

Two more groups complete the map. \emph{Network engineers} own the links from the firm to its venues: the colocated
racks and cross-connects (Book~14, chapter~9), the switches and layer-1 devices (chapter~2), the clocks that timestamp
every packet (chapter~4), the capture that records them (chapter~5), and, for the firms that race between cities, the
long-haul fibre and wireless links (chapters~13 and~14). \emph{Infrastructure engineers} own what research and trading
run on away from the exchange: the compute clusters and their schedulers, storage for tick data, the build and release
systems, and the cloud accounts that some firms now use for research (Book~15). The two groups are measured by capacity,
cost and failure rather than by microseconds, but a failure in either stops trading as surely as a bug in the strategy.

\begin{definition}[Site reliability engineer]\label{def:in:software-engineer:sre}
A \emph{site reliability engineer}\index{site reliability engineer} is a software engineer who owns the reliability of
production systems: monitoring and alerting, the on-call rotation, incident response and review, capacity, and the
automation that removes manual operational work.
\end{definition}

The name comes from Google, whose book on the practice describes it as ``what happens when you ask a software engineer to
design an operations team'' and caps operational work at half of the team's time so that the other half goes to
engineering it away. Trading firms often call the job production engineering or trading support; the tools are those of
Book~15, chapter~28: service-level objectives, the on-call rotation and blameless reviews.

\section{How a trading firm differs from a technology company}

Four differences shape an engineer's day.
\begin{itemize}
\item \textbf{Distance to the money.} A trading firm's system earns or loses money directly, within seconds of a change.
An engineer can see the result of a release in the next day's profit and loss, and a failure can cost more in minutes
than the engineer's year's pay. A technology company's engineer is usually several layers from revenue.
\item \textbf{The market's clock.} A trading system must be up while its markets are open and can change while they are
closed. Releases go out between sessions, on-call matters most while markets are open, and weekends are for maintenance.
Crypto markets, which never close, remove that rhythm (chapter~26).
\item \textbf{Size.} A trading firm of a thousand people has a few hundred engineers; the product teams and platform teams
of Book~16, chapter~19, have a few people each, and one engineer may own a system end to end.
\item \textbf{What cannot be bought.} A trading firm builds some of what other companies buy as a service, when latency,
control or secrecy requires it (Book~16, chapter~20); and much of what it builds, it does not publish.
\end{itemize}

\begin{example}[An error budget on the market's clock]\label{ex:in:software-engineer:budget}
A service-level objective of 99.9\% availability measured over the whole year allows 525.6 minutes of failure. Measured
over US equity trading hours alone (6.5 hours on 252 days, 1\,638 hours), the same 99.9\% allows 98.3 minutes, and 99.99\%
allows 9.8. The trading firm's budget is smaller, and every minute of it falls when the market is moving.
\end{example}

\section{Pay and progression compared with technology firms}

The pay evidence comes from two sources that answer different questions. The labour condition applications of
chapter~14 compare trading firms with banks and exchanges for the same occupation code at the same \omterm{def:in:pay-levels-by-role-firm-type-and-seniority:soc}{wage level}, with
intervals; they do not cover technology employers here, because the one pass over the files kept only the sourced finance
employers. The occupational survey compares industries, including technology, but only as published percentiles without
intervals.

\begin{dated}{2025-09}{What software engineers are paid: the public evidence}\label{dat:in:software-engineer:pay}
Fiscal 2025 labour condition applications coded as software developers (15-1252), median offered base: systematic funds
\$207\,000 (67 applications, 6 employers), platforms \$200\,000 (35, 3), market makers \$175\,000 (153, 9), banks
\$155\,688 (4\,790, 8), exchanges \$147\,900 (296, 4). Trading firms over banks, ratio of medians 1.19 (95\% interval
1.16--1.23), 1.17 at level IV. Occupational survey, May 2025, software developers' median: web search portals and other
information \$213\,520; software publishers \$164\,550; securities industry \$163\,290; computing infrastructure and data
processing \$156\,950; banks \$136\,620; computer systems design \$132\,050. \omterm{def:in:how-pay-works:base}{Base salary} only.
\end{dated}

\begin{omfigure}
\begin{tikzpicture}
\begin{axis}[width=10.5cm,height=5.2cm,xmin=80,xmax=320,ymin=0.4,ymax=5.6,ytick=data,
  yticklabels from table={figdata/industry/20-software-engineer/kinds.csv}{kind},yticklabel style={font=\footnotesize},
  xticklabel style={font=\footnotesize},xlabel={\small offered base, \$ thousand},grid=major,grid style={gray!20},
  table/col sep=comma]
\addplot[draw=none,mark=none,error bars/.cd,x dir=both,x explicit,error mark=none,error bar style={omDef,line width=1pt}]
  table[x=p50,y=pos,x error plus expr=\thisrow{p90}-\thisrow{p50},x error minus expr=\thisrow{p50}-\thisrow{p10}]
  {figdata/industry/20-software-engineer/kinds.csv};
\addplot[draw=none,mark=none,error bars/.cd,x dir=both,x explicit,error mark=none,error bar style={omDef!40,line width=7pt}]
  table[x=p50,y=pos,x error plus expr=\thisrow{p75}-\thisrow{p50},x error minus expr=\thisrow{p50}-\thisrow{p25}]
  {figdata/industry/20-software-engineer/kinds.csv};
\addplot[only marks,mark=|,mark size=5pt,line width=1.5pt,omThm] table[x=p50,y=pos]{figdata/industry/20-software-engineer/kinds.csv};
\end{axis}
\end{tikzpicture}

\omcaption{Software developers' offered base in fiscal 2025 labour condition applications by kind of employer (code
15-1252, all titles): 10th to 90th percentile (thin), interquartile range (thick), median (mark). Data:
\texttt{data/industry/lca\_soc.csv}, through \texttt{in\_engineer.cells}.}\label{fig:in:software-engineer:kinds}
\end{omfigure}

The market makers' range is the widest: their 90th percentile, \$295\,000, is the highest of any kind, and their 10th is
close to the banks'. Pooled and compared level by level, the trading firms' premium can be put in a single number with an
interval.

\begin{omfigure}
\begin{tikzpicture}
\begin{axis}[width=10.5cm,height=5.2cm,xmin=0.5,xmax=5.5,ymin=0.7,ymax=2.3,xtick={1,2,3,4,5},
  xticklabels={all levels,I,II,III,IV},xlabel={\small wage level},ylabel={\small ratio of medians},
  tick label style={font=\footnotesize},grid=major,grid style={gray!20},table/col sep=comma,unbounded coords=jump,
  legend style={legend cell align=left,font=\scriptsize,at={(0.5,-0.3)},anchor=north,/tikz/every even column/.append style={column sep=8pt}},
  legend columns=3]
\addplot[only marks,mark=o,black!60,thick,error bars/.cd,y dir=both,y explicit]
  table[x expr=\thisrow{pos}-0.15,y=bank2021,y error plus expr=\thisrow{bank2021_hi}-\thisrow{bank2021},
  y error minus expr=\thisrow{bank2021}-\thisrow{bank2021_lo}]{figdata/industry/20-software-engineer/ratio.csv};
\addplot[only marks,mark=*,omDef,thick,error bars/.cd,y dir=both,y explicit]
  table[x=pos,y=bank2025,y error plus expr=\thisrow{bank2025_hi}-\thisrow{bank2025},
  y error minus expr=\thisrow{bank2025}-\thisrow{bank2025_lo}]{figdata/industry/20-software-engineer/ratio.csv};
\addplot[only marks,mark=square*,omThm,thick,error bars/.cd,y dir=both,y explicit]
  table[x expr=\thisrow{pos}+0.15,y=exch2025,y error plus expr=\thisrow{exch2025_hi}-\thisrow{exch2025},
  y error minus expr=\thisrow{exch2025}-\thisrow{exch2025_lo}]{figdata/industry/20-software-engineer/ratio.csv};
\draw[black!50] (axis cs:0.5,1) -- (axis cs:5.5,1);
\legend{over banks 2021,over banks 2025,over exchanges 2025}
\end{axis}
\end{tikzpicture}

\omcaption{Median offered base of software-developer filings at trading firms (market makers, systematic funds,
platforms) over banks' and exchanges', overall and by \omterm{def:in:pay-levels-by-role-firm-type-and-seniority:soc}{wage level}, with 95\% bootstrap intervals. Missing points are
comparisons with fewer than ten filings or three employers on a side. Data: \texttt{data/industry/lca\_kind\_gap.csv},
through \texttt{in\_engineer.kind\_gaps}.}\label{fig:in:software-engineer:ratio}
\end{omfigure}

The trading firms' premium over banks is about a fifth at every level where it can be measured in fiscal 2025, and it
narrowed from about a third in fiscal 2021 as the banks' offers rose faster (\cref{fig:in:software-engineer:ratio}). Over
exchanges it is larger. \omterm{def:in:how-pay-works:base}{Base salary} understates the difference: chapter~14's implied ratios of total pay to base were
several times higher at market makers than at banks.

\begin{omfigure}
\begin{tikzpicture}
\begin{axis}[width=10.5cm,height=5.6cm,xmin=0.7,xmax=5.3,ymin=60,ymax=300,xtick={1,2,3,4,5},
  xticklabels={10th,25th,median,75th,90th},xlabel={\small percentile},ylabel={\small annual wage, \$ thousand},
  tick label style={font=\footnotesize},grid=major,grid style={gray!20},table/col sep=comma,
  legend style={font=\scriptsize,at={(0.5,-0.3)},anchor=north,legend cell align=left,/tikz/every even column/.append style={column sep=6pt}},
  legend columns=3]
\addplot[omThm,thick,mark=square*] table[x=pos,y=search]{figdata/industry/20-software-engineer/survey.csv};
\addplot[omDef,thick,mark=*] table[x=pos,y=securities]{figdata/industry/20-software-engineer/survey.csv};
\addplot[black!70,thick,mark=triangle*] table[x=pos,y=publishers]{figdata/industry/20-software-engineer/survey.csv};
\addplot[omDef!50,thick,dashed,mark=o,mark options={solid}] table[x=pos,y=banks]{figdata/industry/20-software-engineer/survey.csv};
\addplot[black!40,thick,dashed,mark=diamond*,mark options={solid}] table[x=pos,y=systems]{figdata/industry/20-software-engineer/survey.csv};
\legend{web search and information,securities,software publishers,banks,computer systems design}
\end{axis}
\end{tikzpicture}

\omcaption{Software developers' annual wages by industry at five percentiles, May 2025 occupational survey. The securities
industry pays like the software publishers at the median and above them at both ends; web search and information pays more
at every percentile. Data: \texttt{data/industry/oews\_roles.csv}, through \texttt{in\_engineer.survey}.}\label{fig:in:software-engineer:survey}
\end{omfigure}

The survey places the securities industry beside the software publishers, not above them: the ratio of their percentiles
is 1.27 at the 10th, 0.99 at the median and 1.12 at the 90th (\cref{fig:in:software-engineer:survey}). The industry that
contains the largest internet companies pays more at every percentile: the securities industry's median is 0.77 of
theirs. Two cautions apply. The securities industry in the survey includes brokers and asset managers as well as the
trading firms, whose filings sit above the survey's median; and the survey measures wages, not the equity awards that make
up much of a technology engineer's pay (chapter~13). The honest conclusion is narrow: for base pay, a trading firm competes
with the technology industry's best-paying part, not above it.

Progression differs with size. An employer with tens of thousands of engineers needs a ladder of levels to compare them
with one another; a firm with a few hundred can judge each engineer by the systems they own, and its pay follows the
firm's year (chapter~13). Chapter~28 follows the paths.

\section{Public engineering: blogs, talks and open source}

Most of a trading firm's code is private, but some firms publish a great deal, and what they publish is the best public
evidence of what their engineers do.

\begin{dated}{2026-09}{Trading firms' public engineering}\label{dat:in:software-engineer:public}
Jane Street writes ``everything that we can in OCaml'', which it calls its ``primary development platform'', publishes
open-source OCaml libraries and tools (Base, the Dune build system, the OxCaml language extensions, and Hardcaml, ``A
language for designing and simulating hardware in OCaml''), has funded work on the OCaml compiler and its package manager,
and runs a public engineering blog; its GitHub organisation listed 411 public repositories on 29 September 2026. Jump
Crypto's Firedancer, a validator client for the Solana blockchain, is open source; its README says that ``The concurrency
model draws from experience in the low latency trading space''.
\end{dated}

For a candidate, public engineering answers questions that interviews do not: which language the firm writes in, how it
tests, what it considers a hard problem. For the firm it is recruiting, and a way to have outside engineers improve the
tools it depends on. What is never published is the trading logic, and the engineer who joins accepts that most of their
work will stay private (chapter~28 discusses what a candidate can show from such a job).

\section{Tutorial: trading firms, banks and technology}\label{tut:in:software-engineer:tut}

\begin{tutorial}
\textbf{Goal.} Compare software engineers' pay at trading firms with banks, exchanges and technology industries, with the
uncertainty each source allows.
\textbf{End state:} \cref{fig:in:software-engineer:kinds,fig:in:software-engineer:ratio,fig:in:software-engineer:survey}.
\begin{tutsteps}
\item \textbf{Cells by kind.} Chapter~19's \texttt{lca\_soc.csv} for code 15-1252 by employer kind, at all levels.
\item \textbf{Ratio of medians.} \texttt{in\_lca\_kind\_gap\_derive.py} pools the trading firms and applies
\texttt{firm.roles.median\_gap} against banks and against exchanges, overall and by level, from chapter~14's cache,
keeping a comparison only when each side has ten filings from three employers.
\item \textbf{Percentile ratios.} \texttt{percentile\_ratio(a, b)} divides two published distributions percentile by
percentile and returns nothing where a value is marked (\cref{lst:in:software-engineer:ratio}).
\omcode{code/firm/roles/firm_roles.py}{284}{293}{Two published wage distributions compared percentile by percentile.}\label{lst:in:software-engineer:ratio}
\item \textbf{Scale.} The headcounts from the filings give the hook's ratios.
\end{tutsteps}
\end{tutorial}

For fiscal 2025, trading firms over banks: 1.19 (1.16--1.23) overall, 1.25 (1.14--1.38) at level~II, 1.18 (1.12--1.27) at
level~III, 1.17 (1.11--1.29) at level~IV; level~I's interval (0.78--1.43) includes one. Over exchanges: 1.25 overall and
1.30 at level~IV. Securities over software publishers in the survey: 1.27, 1.07, 0.99, 1.00, 1.12 from the 10th to the
90th percentile.

\textbf{What to change next.} Add large technology employers to the filings with a sourced list, which needs a second
pass over the files; weight the comparison by level so that the overall ratio is not a mix effect (chapter~19); add bonus
and equity from the proxy statements' \omterm{def:in:pay-levels-by-role-firm-type-and-seniority:median}{median-employee pay} (chapter~14).

\section{Build: engineering cards and percentile comparison}\label{bld:in:software-engineer:roles}

\begin{build}
\bfield{Purpose} Add the three engineering specialisms to the role registry and a comparison of published distributions.
\bfield{Interface} \texttt{firm.roles}: the cards \texttt{low-latency engineer}, \texttt{FPGA engineer},
\texttt{site reliability engineer}; \texttt{percentile\_ratio(a, b, keys) -> \{key: ratio or None\}}; with
\texttt{median\_gap} (chapter~19) for data with values.
\bfield{Rules} A marked or missing percentile gives no ratio, never a guess; ratios of published percentiles carry no
interval and the chapter says so.
\bfield{Acceptance tests} \texttt{code/firm/roles/tests/}: known ratios; a top-code mark gives \texttt{None}.
\bfield{Stretch} Bounds for a ratio when one side is top-coded; interpolate a distribution from its five percentiles and
compare shares above a threshold.
\end{build}

\begin{omsources}
\item Virtu Financial, Form 10-K for 2025; Flow Traders, Annual Report 2025; Microsoft, Form 10-K for fiscal 2026.
\item Beyer, B., C.~Jones, J.~Petoff and N.~R.~Murphy (eds., 2016), \emph{Site Reliability Engineering}, O'Reilly.
\item Jane Street technology page and GitHub organisation; Firedancer repository.
\item Chapter~14's tables from the Department of Labor's LCA files; BLS occupational survey, May 2025.
\end{omsources}

\section{Exercises}

\begin{exercise}[$\star$]\label{exo:in:software-engineer:1}
How many times Virtu's staff is Microsoft's product research and development staff, and its whole staff?
\end{exercise}

\begin{exercise}[$\star$]\label{exo:in:software-engineer:2}
How many minutes of failure a year does a 99.95\% objective allow over US equity trading hours?
\end{exercise}

\begin{exercise}[$\star$]\label{exo:in:software-engineer:3}
Flow Traders' 2025 annual report gives its 2024 year-end staff as 609 FTEs in one section and 625 in another. What should
a reader do with such a difference?
\end{exercise}

\begin{exercise}[$\star\star$]\label{exo:in:software-engineer:4}
Why is the trading firms' premium over exchanges larger at level~II than at level~IV?
\end{exercise}

\begin{exercise}[$\star\star$]\label{exo:in:software-engineer:5}
The ratio of trading firms' to banks' median fell from 1.32 in fiscal 2021 to 1.19 in fiscal 2025. Using the medians of
the derived table (\$165\,000 and \$125\,000 in 2021; \$185\,000 and \$155\,688 in 2025), which side moved more?
\end{exercise}

\begin{exercise}[$\star\star$]\label{exo:in:software-engineer:6}
An engineer choosing between a trading firm and a large technology company compares base salaries only. What does that
leave out on each side?
\end{exercise}

\begin{exercise}[$\star\star\star$]\label{exo:in:software-engineer:7}
\emph{Coding.} With \texttt{percentile\_ratio}, compare the securities industry's software developers with each of the
other five industries of the survey table. At how many of the 25 percentile comparisons does the securities industry pay
more?
\end{exercise}

\begin{exercise}[$\star\star\star$]\label{exo:in:software-engineer:8}
\emph{Find the flaw.} ``The survey shows the securities industry pays software developers no more than software
publishers do, so trading firms pay engineers like the technology industry.''
\end{exercise}

\section{Problem: Trading Firm or Technology Company?}

\begin{problem}[{Weekend problem --- trading firm or technology company?}]\label{pb:in:software-engineer:1}
A software engineer with five years at a large technology company has an offer from a market maker. She wants to know how
the work and the pay differ.

\medskip\noindent\textbf{Part I --- The jobs.}
\begin{enumerate}
\item Define the \omterm{def:in:software-engineer:lowlat}{low-latency engineer}, the \omterm{def:in:software-engineer:fpga}{FPGA engineer} and the \omterm{def:in:software-engineer:sre}{site reliability engineer}.
\item What do core engineers build?
\item Where does the \omterm{def:in:software-engineer:sre}{site reliability engineer}'s name come from, and what rule on operational work goes with it?
\item What makes the \omterm{def:in:software-engineer:fpga}{FPGA engineer}'s testing stricter than a software engineer's?
\item Which books of the series teach each specialism?
\end{enumerate}

\medskip\noindent\textbf{Part II --- The firm.}
\begin{enumerate}[resume]
\item Compare the scale of the two employers with the hook's numbers.
\item State the four differences between a trading firm and a technology company.
\item Compute the error budget of a 99.9\% objective over a year and over trading hours.
\item What does the market's clock do to releases and on-call?
\item What do trading firms publish about their engineering, and what not?
\end{enumerate}

\medskip\noindent\textbf{Part III --- The pay.}
\begin{enumerate}[resume]
\item What can each source (filings, survey) compare, and what not?
\item Give the trading firms' medians by kind and the banks'.
\item Give the ratios over banks by level, with intervals.
\item Compare the securities industry with software publishers and with web search in the survey.
\item What does \omterm{def:in:how-pay-works:base}{base salary} leave out on each side?
\end{enumerate}

\medskip\noindent\textbf{Part IV --- The verdict.}
\begin{enumerate}[resume]
\item State the \emph{named result}: the ratio of median offered base at trading firms to banks for software developers
at each \omterm{def:in:pay-levels-by-role-firm-type-and-seniority:soc}{wage level}, with bootstrap intervals, and the survey's comparison with technology industries.
\item Why is the comparison with banks sharper than the comparison with technology?
\item How does progression differ?
\item What should she ask the market maker about on-call?
\item In two sentences, answer her.
\end{enumerate}
\end{problem}

\section{Interview questions}

\begin{interviewq}[$\star$ \iqroles{developer}]\label{iq:in:software-engineer:1}
What is tail latency, and why does a trading firm care more about the 99.9th percentile than the mean?
\end{interviewq}

\begin{interviewq}[$\star$ \iqroles{developer}]\label{iq:in:software-engineer:2}
A service must be up 99.9\% of trading hours. How much downtime is that a year, and how would you spend it?
\end{interviewq}

\begin{interviewq}[$\star\star$ \iqroles{developer}]\label{iq:in:software-engineer:3}
How does kernel bypass reduce latency, and what do you give up?
\end{interviewq}

\begin{interviewq}[$\star\star$ \iqroles{developer, trader}]\label{iq:in:software-engineer:4}
Your system sent duplicate orders for two seconds at the open. Walk through the first hour after.
\end{interviewq}

\begin{interviewq}[$\star\star$ \iqroles{developer}]\label{iq:in:software-engineer:5}
When would you put a risk check in an FPGA rather than in software?
\end{interviewq}

\begin{interviewq}[$\star\star\star$ \iqroles{developer}]\label{iq:in:software-engineer:6}
Design the release process for a trading system that must not change while the market is open but must ship every day.
\end{interviewq}
